\documentclass[a4paper,11pt]{article}
\usepackage[T1]{fontenc}
\usepackage[utf8]{inputenc}
\usepackage[ngerman]{babel}
\usepackage{lmodern}
\usepackage{amsmath,amssymb}
\usepackage{geometry}
\usepackage[table]{xcolor}
\usepackage{tikz}
\usepackage{eso-pic}
\geometry{paper=a4paper,top=0cm,bottom=0cm,left=0cm,right=0cm,headheight=0pt,headsep=0pt,footskip=0pt}
\pagestyle{empty}
\definecolor{examblue}{RGB}{0,70,130}
\definecolor{examgreen}{RGB}{0,110,60}
\definecolor{cardgray}{RGB}{248,248,248}
\newlength{\cardw}\newlength{\cardh}
\setlength{\cardw}{9.80cm}\setlength{\cardh}{5.24cm}
\AddToShipoutPictureFG{\begin{tikzpicture}[remember picture,overlay]
\draw[gray!70,densely dotted,line width=0.7pt]([xshift=10.5cm]current page.south west)--([xshift=10.5cm]current page.north west);
\foreach \y in {5.94,11.88,17.82,23.76}{\draw[gray!70,densely dotted,line width=0.7pt]([yshift=\y cm]current page.south west)--([yshift=\y cm]current page.south east);}
\end{tikzpicture}}
\newcommand{\checkfeld}{\raisebox{-0.15ex}{\fbox{\rule{0pt}{1.15ex}\rule{1.15ex}{0pt}}}}
\newcommand{\frontcard}[3]{\begin{tikzpicture}
\node[draw=examblue,line width=0.8pt,rounded corners=4pt,fill=white,minimum width=\cardw,minimum height=\cardh,text width=8.75cm,align=center,inner sep=8pt](card){\Large\bfseries #3};
\node[anchor=north west,fill=examblue,text=white,rounded corners=3pt,font=\bfseries\small,inner xsep=7pt,inner ysep=4pt]at([xshift=7pt,yshift=-7pt]card.north west){Karte #1};
\node[anchor=south,text=examblue,font=\scriptsize\bfseries]at([yshift=21pt]card.south){NaWi $\cdot$ #2};
\node[anchor=south,text=examblue,font=\scriptsize]at([yshift=6pt]card.south){Gewusst:\enspace\checkfeld\enspace\checkfeld\enspace\checkfeld\enspace\checkfeld\enspace\checkfeld};
\end{tikzpicture}}
\newcommand{\backcard}[3]{\begin{tikzpicture}
\node[draw=examgreen,line width=0.8pt,rounded corners=4pt,fill=cardgray,minimum width=\cardw,minimum height=\cardh,text width=8.75cm,align=center,inner sep=9pt](card){\large #3};
\node[anchor=north west,fill=examgreen,text=white,rounded corners=3pt,font=\bfseries\small,inner xsep=7pt,inner ysep=4pt]at([xshift=7pt,yshift=-7pt]card.north west){Karte #1};
\node[anchor=south,text=examgreen,font=\scriptsize\bfseries]at([yshift=21pt]card.south){NaWi $\cdot$ #2};
\node[anchor=south,text=examgreen,font=\scriptsize]at([yshift=6pt]card.south){Gewusst:\enspace\checkfeld\enspace\checkfeld\enspace\checkfeld\enspace\checkfeld\enspace\checkfeld};
\end{tikzpicture}}
\newcommand{\blankcard}{\begin{tikzpicture}\node[minimum width=\cardw,minimum height=\cardh]{};\end{tikzpicture}}
\newcommand{\cardcell}[1]{\parbox[c][5.94cm][c]{10.50cm}{\centering #1}}
\newcommand{\cardrow}[2]{\noindent\cardcell{#1}\cardcell{#2}\par}
\begin{document}
\setlength{\topskip}{0pt}
\offinterlineskip
% Karten 149 bis 159: Vorder- und Rueckseiten
\cardrow
{\frontcard{149}{SRT \& ART}{Womit beschäftigt sich die Relativitätstheorie?}}
{\frontcard{150}{SRT \& ART}{Nenne zumindest drei der fünf wichtigsten Vorhersagen der speziellen Relativitätstheorie.}}
\cardrow
{\frontcard{151}{SRT \& ART}{Was sind die beiden Grundannahmen der SRT?}}
{\frontcard{152}{SRT \& ART}{Erläutere die Zeitdilatation.}}
\cardrow
{\frontcard{153}{SRT \& ART}{Erläutere die Längenkontraktion.}}
{\frontcard{154}{SRT \& ART}{Was versteht man unter der relativistischen Massenzunahme?}}
\cardrow
{\frontcard{155}{SRT \& ART}{Was verstehen wir unter Raumzeit?}}
{\frontcard{156}{SRT \& ART}{Nenne zumindest drei wesentliche Erkenntnisse aus der ART.}}
\cardrow
{\frontcard{157}{SRT \& ART}{Was ist ein Schwarzes Loch?}}
{\frontcard{158}{SRT \& ART}{Was verstehen wir unter gravitativer Zeitdilatation?}}
\newpage
\cardrow
{\backcard{150}{SRT \& ART}{Zum Beispiel Zeitdilatation, Längenkontraktion, Relativität der Gleichzeitigkeit, relativistische Geschwindigkeitsaddition und \(E=mc^2\).}}
{\backcard{149}{SRT \& ART}{Mit Raum, Zeit, Bewegung und Gravitation; SRT für Inertialsysteme, ART für Gravitation als Raumzeitkrümmung.}}
\cardrow
{\backcard{152}{SRT \& ART}{Eine bewegte Uhr geht langsamer: \(\Delta t=\gamma\Delta\tau\).}}
{\backcard{151}{SRT \& ART}{Die Naturgesetze sind in allen Inertialsystemen gleich; die Vakuumlichtgeschwindigkeit ist für alle gleich.}}
\cardrow
{\backcard{154}{SRT \& ART}{Historisch \(m_\mathrm{rel}=\gamma m_0\); heute: Ruhemasse bleibt konstant, Energie und Impuls wachsen.}}
{\backcard{153}{SRT \& ART}{Ein bewegter Körper wird in Bewegungsrichtung kürzer: \(L=L_0/\gamma\).}}
\cardrow
{\backcard{156}{SRT \& ART}{Gravitation ist Raumzeitkrümmung; Lichtablenkung; gravitative Zeitdilatation; Gravitationswellen; Schwarze Löcher.}}
{\backcard{155}{SRT \& ART}{Eine vierdimensionale Struktur aus drei Raumdimensionen und einer Zeitdimension.}}
\cardrow
{\backcard{158}{SRT \& ART}{Uhren laufen in stärkeren Gravitationsfeldern langsamer.}}
{\backcard{157}{SRT \& ART}{Ein Raumzeitbereich, aus dem innerhalb des Ereignishorizonts nichts, auch kein Licht, entkommen kann.}}
\newpage
\cardrow
{\frontcard{159}{SRT \& ART}{Was sind Gravitationswellen?}}
{\blankcard}
\cardrow
{\blankcard}
{\blankcard}
\cardrow
{\blankcard}
{\blankcard}
\cardrow
{\blankcard}
{\blankcard}
\cardrow
{\blankcard}
{\blankcard}
\newpage
\cardrow
{\blankcard}
{\backcard{159}{SRT \& ART}{Wellen der Raumzeitkrümmung, die sich mit Lichtgeschwindigkeit ausbreiten und durch beschleunigte unsymmetrische Massen entstehen.}}
\cardrow
{\blankcard}
{\blankcard}
\cardrow
{\blankcard}
{\blankcard}
\cardrow
{\blankcard}
{\blankcard}
\cardrow
{\blankcard}
{\blankcard}
\end{document}
